\clearpage
\section{A guide to the higher chapters}\label{c:higher}
The comparison theorem is already complete before the higher chapters.
Those chapters pursue a new question: what additional structure is retained
by richer quotients and more informative probes? The following routes use
the \emph{main exposition's} section numbers.

\textbf{First bridge: stay with ordinary algebra.}
\mainref{9}{Main §9} studies attachment data and marked extensions:
besides the two end objects, retain how one sits inside the total object
and how the total object maps onto the other. It also explains why data
from adjacent layers can miss information across a skipped layer.
\mainref{10}{Main §10} studies probes for nilpotent operators and
reconstructs their chains. Modules and linear algebra suffice for this
route. Exercise 7 gives the simplest reason to want it: a factor list
does not determine the assembled object.

\textbf{First higher bridge: an orbit can hide a stabilizer.}
Start with \mainref{12-2}{main §12.2}. A groupoid is a category in which
every arrow is invertible. For a group action, its action groupoid has
the acted-on points as objects and the group elements carrying one point
to another as arrows.

Let \(G=C_6\) act by addition on the two-element set \(G/B\).
The action is transitive, so its ordinary orbit set is a single point.
But the stabilizer of either point is \(B=C_3\). The action groupoid
retains these automorphisms. Its full subgroupoid on the object \(B\)
has one object and automorphism group \(B\); it is equivalent to the
whole groupoid because every other object is isomorphic to \(B\).
Writing \(BK\) for the classifying space of a group \(K\), we get
\[
 (G/B)/G=*,\qquad (G/B)//G\simeq BB.
\]
The double slash denotes the homotopy quotient, obtained by realizing
the action groupoid. It remembers something the orbit set discarded.

\textbf{Next bridge: intersections acquire components and loops.}
Read \mainref{12-3}{main §12.3} for the homotopy-pullback construction.
The double-coset formula proved in \mainref{12-6}{main §12.6} gives,
for subgroups \(H,K\leq G\),
\[
 BH\times^h_{BG}BK
 \simeq\coprod_{[g]\in H\backslash G/K} B(H\cap gKg^{-1}).
\]
Here the superscript \(h\) denotes a homotopy pullback. In our abelian
example, \(A+B=G\) and \(A\cap B=0\), so
\[
 BA\times^h_{BG}BB\simeq *.
\]
With \(A\) on both sides there are three double cosets and each
intersection is \(A\), giving
\[
 BA\times^h_{BG}BA\simeq BA\sqcup BA\sqcup BA.
\]
An ordinary subgroup intersection sees the group in the identity
component. The homotopy pullback records the other components too.

\pausehere{At each stage ask which object is being used. The orbit set
forgets stabilizers; the action groupoid and its classifying homotopy
\(1\)-type retain equivalent information.}

\clearpage
\textbf{Choose a route by the question you want to answer.}
The higher chapters need not be read in numerical order.

\begin{center}
\small
\begin{tabular}{@{}>{\raggedright\arraybackslash}p{0.17\linewidth}>{\raggedright\arraybackslash}p{0.31\linewidth}>{\raggedright\arraybackslash}p{0.43\linewidth}@{}}
\toprule
\textbf{Main sections}&\textbf{Preparation}&\textbf{Question and payoff}\\
\midrule
\mainref{12-1}{§12.1},
\mainref{12-4}{§§12.4--5},
\mainref{12-7}{§12.7}
&Spaces, homotopy pullbacks and universal properties
&How do coherent algebraic models, base change and descent organize
quotients? What extra structure replaces ordinary normality?\\[5pt]
\mainref{13-1}{§§13.1--2}
&Groups, semidirect products and modules
&What does linearization retain when performed over a fixed group?
The base retains the group action on coefficients.\\[5pt]
\mainref{13-3}{§13.3}
&Spectra, truncations and descent
&How do spectra over \(BG\) extend these coefficients across degrees?
The heart recovers ordinary \(G\)-modules.\\[5pt]
\mainref{13-4}{§§13.4--5}
&Complexes, cohomology, projective resolutions and derived categories
&How are marked extensions and their automorphisms recorded?
Why can derived tensor products have additional homology?\\[5pt]
\mainref{13-6}{§13.6}
&Derivations and commutative algebra, followed by cotangent complexes
&How does cotangent transitivity identify the obstruction to extending
a derivation?\\
\bottomrule
\end{tabular}
\end{center}

For an algebraic entry to this table, read \mainref{13-1}{§13.1} and
\mainref{13-2}{§13.2} before
attempting spectra. An abelian group object over \(G\) is equivalent to a
\(\Z[G]\)-module:
\[
 \operatorname{Ab}(\operatorname{Grp}/G)\simeq\Z[G]\text{-}\operatorname{Mod}.
\]
The zero section splits the map to \(G\); its kernel is abelian, and
conjugation by that section gives the \(G\)-action. Conversely a
\(G\)-module produces its semidirect product over \(G\).
Thus linearizing \emph{over a base} retains more than forgetting the base
and taking an ordinary abelian group.

\textbf{Three questions keep the higher route precise.}
First, what is the ambient category? A descent result for an
\(\infty\)-topos does not automatically hold in every category of
algebraic models. Second, what extra data are retained---a base, marked
endpoints, automorphisms, or coherent homotopies? Third, what is actually
concluded? An obstruction group names possible obstructions; a particular
lifting problem has a particular obstruction class whose vanishing must
be checked.

The main exposition identifies the higher foundations it imports.
These chapters enrich quotient and reconstruction questions; they do
not claim a Jordan--Hölder theorem for arbitrary higher objects.
The ordinary refinement proof does not depend on them.

\route{Further reading: \href{https://tensorgami.com/notes/quotients-and-reconstruction.html}
{Quotients, Refinement and Reconstruction};
\href{https://tensorgami.com/notes/quotients-and-reconstruction-audit.html}
{notes on proofs, dependencies, and revisions}.}
